% appendix-notation.tex — \input'd from main.tex (after \appendix)
%
% Sources: Jörn's definitions dictated in chat. No other source.
%
\section{Notation Glossary}\label{app:notation}

% ---------- 1. Notation table ----------


\begin{longtable}{p{5.3cm} p{3.4cm} p{3.4cm}}
\toprule
\textbf{Concept} & \textbf{This thesis} & \textbf{HK2017} \\
\midrule
\endfirsthead
\midrule
\textbf{Concept} & \textbf{This thesis} & \textbf{HK2017} \\
\midrule
\endhead

Ambient space                                    & $\mathbb{R}^4$      & $\mathbb{R}^{2n}$ \\
Coordinates                                      & $(q,p)$              & --- \\
Std.\ complex structure                          & $J_0$                & $J$ \\
Std.\ symplectic form                            & $\omega_0$           & $\omega$ \\
Liouville form                                   & $\lambda_0$          & --- \\
Convex body                                      & $K$                  & $K$ \\
Facet count                                      & $F$                  & $\mathbf{F}_K$ \\
Facets                                           & $F_i$                & $F_i$ \\
Outward unit normals                             & $n_i$                & $n_i$ \\
Heights                                          & $h_i$                & $h_i$ \\
Support function                                 & $h_K$                & $h_K$ \\
Gauge function                                   & $g_K$                & $g_K$ \\
Hamiltonian                                      & $g_K^2$              & $g_K^2$ \\
Normal cone                                      & $N_K(x)$             & $N_K(x)$ \\
Reeb vectors (per-facet)                         & $R_i$                & $p_i$ \\
Curve domain                                     & $[0,T]$              & $[0,1]$ \\
Gen.\ closed char.\ / Reeb orbit / Ham.\ orbit  & $\gamma$             & $\gamma$ \\
Action                                           & $A(\gamma)$          & $A(\gamma)$ \\
EHZ capacity                                     & $c_\mathrm{EHZ}(K)$ & $c_\mathrm{EHZ}(K)$ \\
Dual variable                                    & $z$                  & $z$ \\
Dual functional                                  & $I_K(z)$             & $I_K(z)$ \\
Constraint set                                   & $M(K)$               & $M(K)$ \\
Action matrix                                    & $H$                  & --- \\
Adjacency graph                                  & ---                  & --- \\

\bottomrule
\end{longtable}

A dash (---) indicates the concept is not explicitly named in that source.

% ---------- 2. Notes ----------


\begin{description}

\item[Two normalizations.]
  Ours is \emph{period normalization};
  theirs is \emph{HK2017 normalization}.

\item[MATLAB sign convention.]
  HK2017's MATLAB implementation minimizes~$-Q(\sigma,\beta)$; our Rust code maximizes~$Q$ directly. The optima coincide.

\item[Permutation orientation.]
  Lemma~\ref{lem:H-quadratic} writes the action sum with $j < i$ in
  $\omega_0(n_{\sigma(j)}, n_{\sigma(i)})$.
  The Rust function \texttt{q\_from\_beta} uses $i > j$ (higher index first),
  which negates $Q(\beta)$ by antisymmetry of $\omega_0$.
  Since the code maximises $Q > 0$, it finds the reverse-traversal
  representative of the orbit.
  The action $A = 1/(2Q)$ is unchanged, but
  \texttt{EhzResult::\allowbreak best\_\allowbreak permutation}
  is the cyclic reverse of the orbit described in the thesis.

\end{description}
